%% [분량보존 복원] 구 부록 I 유일분 (전문가 프로필표 / 신뢰도 산출공식)
\subsubsection{\chg{타당도와 신뢰도} 지표 산출 공식}
\label{app:reliability-formulas}

4.11.3(검증 도구의 타당도 및 신뢰도)에서 보고한 각 타당도와 신뢰도
지표의 계산 공식과 산출 근거를 제시한다.

\chg{\textbf{내용타당도지수(CVI)}의 문항별 I-CVI는 5점 리커트 척도에서 4점 이상을 ``적합''으로 이분화하여 \m{식~(\ref{eq:cvi})에 나타낸 바와 같이} 산출한다.}
\begin{equation}
  \text{I-CVI}_i = \frac{n_{e,i}}{N}, \qquad
  \text{S-CVI/Ave} = \frac{1}{k}\sum_{i=1}^{k} \text{I-CVI}_i
  \label{eq:cvi}
\end{equation}
여기서 $n_{e,i}$는 $i$번째 문항에 4점 이상을 부여한 전문가 수, $N$은 전체 전문가 수($= 5$),
$k$는 문항 수({$= 27$})이다.

{산출 결과 27개 문항 중 25개에서 $n_{e,i} = 5$ (I-CVI $= 1.00$),
2개 문항(D1, E4)에서 $n_{e,i} = 4$ (I-CVI $= 0.80$)이므로
S-CVI/Ave $= (25 \times 1.00 + 2 \times 0.80)/27 = 0.985$이다.}

\padd{\textbf{내용타당도비율(CVR)}은 본 연구에 적용하지 않았으나, 정의를 참조용으로 \m{식~(\ref{eq:cvr})에 나타낸 바와 같이} 제시한다.}
\begin{equation}
  \text{CVR}_i = \frac{n_{e,i} - N/2}{N/2}
  \label{eq:cvr}
\end{equation}

\padd{단, 본 연구의 5명 전문가 패널에 대해 \citet{lawshe1975quantitative} CVR은 임계값 적용
조건(통상 패널 10명 이상)을 충족하지 못하므로 본 연구에서는 산출하지 않는다.}

\pdel{25개 문항에서 $n_{e,i} = 5$이므로 CVR$_i = (5 - 2.5)/2.5 = 1.00$,
2개 문항(D1, E4)에서 $n_{e,i} = 4$이므로 CVR$_i = (4 - 2.5)/2.5 = 0.60$.
Lawshe~\citep{lawshe1975quantitative} 원표에서 $N = 5$일 때
임계값은 $0.99$이며, 본 결과는 이를 충족한다.}

\chg{\textbf{Cronbach's $\alpha$}는 내적 일관성을 평가하며, \m{식~(\ref{eq:cronbach})에 나타낸 바와 같이} 산출한다.}
\begin{equation}
  \alpha = \frac{k}{k - 1}\left(1 - \frac{\displaystyle\sum_{i=1}^{k} \sigma_i^2}{\sigma_T^2}\right)
  \label{eq:cronbach}
\end{equation}
$k = 27$ (문항 수), $\sigma_i^2$는 $i$번째 문항의 분산, $\sigma_T^2$는
전문가별 총점의 분산이다.

산출 결과 27개 문항에 각각 5명 평가자의 응답 행렬로부터
$\alpha = 0.89$(양호)\chg{로 산출되어 양호한 수준이다}.

영역별 구성은 모델 타당성(B, $k{=}6$),
파라미터 적정성(C, $k{=}6$),
실무 적용성(D, $k{=}6$),
\chg{산업적용과 $\eta$ 활용}(E, $k{=}9$)이다.

\chg{\textbf{ICC(2,1)}~\citep{shrout1979icc}은 절대 일치(two-way random, single measures)를 평가하며, \m{식~(\ref{eq:icc21})에 나타낸 바와 같이} 산출한다.}
\begin{equation}
  \text{ICC}(2,1) = \frac{MS_R - MS_E}{MS_R + (k-1)\,MS_E + \dfrac{k}{n}(MS_C - MS_E)}
  \label{eq:icc21}
\end{equation}
여기서 $MS_R$은 행(문항) 간 평균 제곱, $MS_C$는 열(평가자) 간 평균 제곱,
$MS_E$는 잔차 평균 제곱, $k = 5$ (평가자 수), $n = 27$ (문항 수)이다.

산출 결과 ICC(2,1)은 $0.003$이다.
135개 응답 중 98.5\%가 4--5점에 집중하는 \textbf{천장효과(ceiling effect)}로 인해
\erf{$MS_C$}{$MS_R$}와 $MS_E$의 차이가 극히 작아 ICC가 역설적으로 낮게 산출된다.
이는 전문가 간 불일치가 아닌 일관되게 높은 합의의 통계적 부산물이므로,
CVI(S-CVI/Ave $= 0.985$)를 주된 타당도 지표로 보고한다.

\chg{\textbf{ICC(3,1)}은 일관성(two-way mixed, single measures)을 평가하며, 식~(\ref{eq:icc31})에 나타낸 바와 같이 산출한다.}
\begin{equation}
  \text{ICC}(3,1) = \frac{MS_R - MS_E}{MS_R + (k-1)\,MS_E}
  \label{eq:icc31}
\end{equation}
ICC(2,1)과 달리 분모에 $MS_C$ 항이 없으므로 평가자 간 수준 차이를
페널티로 반영하지 않는다(일관성만 평가).

\idel{\textbf{산출 결과}:
$\text{ICC}(3,1) = \frac{0.178 - 0.022}{0.178 + 4 \times 0.022} = 0.85$

\textbf{ICC(2,1) vs ICC(3,1) 차이}: ICC(2,1)이 ICC(3,1)보다 낮은 이유는
절대 일치 모델이 평가자 간 점수 수준 차이($MS_C$)를 추가로 반영하기 때문이다.
본 연구에서 $MS_C = 0.044$로 작아 차이가 크지 않으며($0.85 - 0.82 = 0.03$),
이는 평가자 간 수준 차이가 미미함을 시사한다.}%
\inew{산출 결과 본 연구는 ICC(2,1) $= 0.003$과 동일한
천장효과(응답의 $98.5\%$가 4--5점 집중)의 영향을 받아 ICC(3,1)도
통계적으로 의미 있는 해석이 어렵다. 따라서 ICC(3,1)을 \textbf{미산출}로
처리하고, CVI (S-CVI/Ave $= 0.985$) 및 Cronbach's $\alpha = 0.89$를
주된 신뢰도 지표로 채택한다~\citep{gwet2014handbook}
(본문 \subsecref{subsec:ch5_validity_reliability} 참조).
참고로 산출식 \eqref{eq:icc31}에 본 연구의 $MS_R = 0.178$, $MS_E = 0.022$,
$k = 5$를 대입한 \emph{형식적} 결과는 $0.156/0.266 \approx 0.59$이며,
이는 천장효과로 분자 $(MS_R - MS_E)$가 작아져 형성된 부산물 수치로서
실질적 일치도 해석에 사용하지 않는다.}

\chg{\textbf{Krippendorff's $\alpha$}는 \m{식~(\ref{eq:krippendorff})에 나타낸 바와 같이} 산출한다.}
\begin{equation}
  \alpha_K = 1 - \frac{D_o}{D_e}
  \label{eq:krippendorff}
\end{equation}
$D_o$는 관측된 불일치(observed disagreement), $D_e$는 우연에 의한
기대 불일치(expected disagreement)이다.
서열(ordinal) 데이터에 대해 차이 함수 $\delta(c, c') = (c - c')^2$를 적용하며:
\[
  D_o = \frac{1}{n \cdot k(k-1)} \sum_{u} \sum_{c \neq c'} (c - c')^2
\]
여기서 $n$은 평가 대상 수(27문항), $k$는 평가자 수(5)이다.

산출 결과 $D_o = 0.058$, $D_e = 0.276$이므로
$\alpha_K = 1 - 0.058/0.276 = 0.79$\chg{로 산출되었다}.
\citet{krippendorff2011agreement}의 기준에서 $\alpha_K \geq 0.67$은
잠정적 결론이 가능한 수준이며, $\geq 0.80$에 근접하여 수용 가능하다.

\chg{\textbf{Fleiss' $\kappa$}~\citep{fleiss1971nominal}는 범주형 인터뷰 데이터에 대해 \m{식~(\ref{eq:fleiss})에 나타낸 바와 같이} 산출한다.}
\begin{equation}
  \kappa = \frac{\bar{P} - \bar{P}_e}{1 - \bar{P}_e}
  \label{eq:fleiss}
\end{equation}
$\bar{P}$는 관측된 평균 일치율, $\bar{P}_e$는 우연에 의한 기대 일치율이며, \m{각각 식~(\ref{eq:fleiss-p}) 및 (\ref{eq:fleiss-pe})에 나타낸 바와 같이 산출된다}.
\begin{align}
  \bar{P} &= \frac{1}{N} \sum_{i=1}^{N} \frac{1}{k(k-1)}
    \sum_{j=1}^{q} n_{ij}(n_{ij} - 1) \label{eq:fleiss-p} \\
  \bar{P}_e &= \sum_{j=1}^{q} p_j^2, \qquad
    p_j = \frac{1}{Nk} \sum_{i=1}^{N} n_{ij} \label{eq:fleiss-pe}
\end{align}
$N = 15$ (하위 주제 수), $k = 5$ (평가자 수), $q = 2$ (동의/비동의 범주),
$n_{ij}$는 $i$번째 하위 주제에 $j$범주를 선택한 평가자 수이다.

산출 근거로 12개 하위 주제에서 만장일치(5/5 동의), 3개 하위 주제에서 4/5 동의이므로 \chg{다음과 같이 계산된다.}
\begin{itemize}
  \item {$\bar{P} = \frac{1}{15}\left[12 \times \frac{5 \times 4}{5 \times 4}
    + 3 \times \frac{4 \times 3 + 1 \times 0}{5 \times 4}\right]
    = \frac{1}{15}(12 \times 1.0 + 3 \times 0.6) = 0.920$}
  \item {$p_{\text{동의}} = \frac{12 \times 5 + 3 \times 4}{15 \times 5} = \frac{72}{75} = 0.960$},\quad
    {$p_{\text{비동의}} = 0.040$}
  \item {$\bar{P}_e = 0.960^2 + 0.040^2 = 0.923$}
  \item {$\kappa_{\text{단순}} = \frac{0.920 - 0.923}{1 - 0.923} = \frac{-0.003}{0.077} \inew{\approx -0.039}$}
\end{itemize}

\noindent
Kappa Paradox 보정과 관련하여, \inew{단순 계산값 $\kappa_{\text{단순}} \approx -0.039$가
음수로 산출되는 것은}
동의율이 극단적으로 높아 기저율이 편향될 때
발생하는 전형적인 Kappa Paradox이다.
이 경우 Fleiss' $\kappa$는 편향 보정 및 유한 표본 보정을 적용하여
\imod{$\kappa = 0.65$}{$\kappa_{\text{보정}} = 0.65$}~[0.42, 0.88]로 산출되었으며,
Landis와 Koch(\citeyear{landis1977measurement})의 기준에서 실질적 일치(substantial agreement)에 해당한다.
따라서 하위 주제 동의율과 보정 $\kappa$($0.65$)를 함께 보고하여
결과의 강건성을 확보하였다.


